\paragraph{事件椭圆的黄金最小伸长律与双曲位移下界}
本段把事件记录 \(\mathbf a=(a_1,\dots,a_T)\) 的矩阵编码 \(M(\mathbf a)\)（定理 \ref{thm:conclusion-elliptic-lossless-lift}）提升为两类可审计的几何量：
其一是椭圆 \(\mathcal E(\mathbf a)=M(\mathbf a)(S^1)\) 的主轴比（条件数），其二是 \(M(\mathbf a)\) 在 \(\mathrm{PSL}_2(\RR)\) 的双曲平移长度。
两者在“固定步数 \(T\)”下都出现严格的黄金最小值，并且极小实现唯一。

\begin{theorem}[固定事件步数的最小伸长：黄金下界与唯一极小词]\label{thm:conclusion-event-ellipse-golden-minimal-stretch}
\leanverified{paper\_conclusion\_event\_ellipse\_golden\_minimal\_stretch}
沿用定理 \ref{thm:conclusion-elliptic-lossless-lift} 的记号。
令 \(\sigma_{\max}(M)\ge \sigma_{\min}(M)>0\) 为矩阵 \(M\in \mathrm{SL}_2(\RR)\) 的奇异值，并定义（\(2\)-范数条件数）
\[
\mathrm{cond}_2(M):=\frac{\sigma_{\max}(M)}{\sigma_{\min}(M)}.
\]
则对任意 \(T\ge 1\) 与任意 \(\mathbf a\in(\NN_{\ge 1})^T\) 都有严格下界
\[
\boxed{\;
\mathrm{cond}_2\!\bigl(M(\mathbf a)\bigr)\ \ge\ \varphi^{4T}\;},
\qquad
\varphi:=\frac{1+\sqrt5}{2},
\]
并且等号成立当且仅当 \(\mathbf a=(1,1,\dots,1)\)。
\leanverified{paper\_event\_ellipse\_golden\_minimal\_stretch\_T1}
\leanverified{R\_mul\_L}
\leanverified{det\_R\_mul\_L}
\leanverified{trace\_R\_mul\_L}
\leanverified{goldenRatio\_sq\_eq}
\leanverified{goldenRatio\_pow\_four\_identity}
\leanverified{goldenRatio\_sq\_mul\_inv\_sq}
\leanverified{goldenRatio\_inv\_eq\_sub\_one}
\leanverified{goldenRatio\_sq\_plus\_inv\_sq}
\leanverified{goldenRatio\_inv\_sq\_char}
\end{theorem}
\begin{proof}
记
\(
R=\bigl(\begin{smallmatrix}1&1\\0&1\end{smallmatrix}\bigr)
\)、
\(
L=\bigl(\begin{smallmatrix}1&0\\1&1\end{smallmatrix}\bigr)
\)，
则对任意 \(m\ge 1\) 有
\[
R^mL=
\begin{pmatrix}
m+1&m\\
1&1
\end{pmatrix}
\succeq
\begin{pmatrix}
2&1\\
1&1
\end{pmatrix}
=RL,
\]
其中 \(\succeq\) 表示逐项不小。
于是
\(
M(\mathbf a)=\prod_{t=1}^T(R^{a_t}L)\succeq (RL)^T
\)。
由于 \(RL\) 为正矩阵，Perron--Frobenius 理论给出谱半径的单调性：
\[
\rho\!\bigl(M(\mathbf a)\bigr)\ \ge\ \rho\!\bigl((RL)^T\bigr)=\rho(RL)^T.
\]
直接计算
\(
RL=\bigl(\begin{smallmatrix}2&1\\1&1\end{smallmatrix}\bigr)
\)
的特征值为
\(
\lambda_\pm=\frac{3\pm\sqrt5}{2}=\varphi^{\pm 2}
\)，
故 \(\rho(RL)=\varphi^2\)，从而 \(\rho(M(\mathbf a))\ge \varphi^{2T}\)。

\par
另一方面，对任意矩阵 \(M\) 有 \(\rho(M)\le \sigma_{\max}(M)\)（谱半径不超过算子范数）。
又因 \(M(\mathbf a)\in \mathrm{SL}_2(\RR)\) 满足 \(\det M(\mathbf a)=1\)，故
\(
\sigma_{\min}(M(\mathbf a))=\sigma_{\max}(M(\mathbf a))^{-1}
\)，
从而
\[
\mathrm{cond}_2\!\bigl(M(\mathbf a)\bigr)
=\sigma_{\max}(M(\mathbf a))^2
\ge \rho(M(\mathbf a))^2
\ge \varphi^{4T}.
\]

\par
若存在某个 \(t\) 使 \(a_t>1\)，则 \(R^{a_t}L\succ RL\)（至少一项严格大于），从而 \(M(\mathbf a)\succ (RL)^T\)。
由于 \((RL)^T\) 为正矩阵，谱半径的严格单调性给出 \(\rho(M(\mathbf a))>\rho((RL)^T)\)，于是 \(\mathrm{cond}_2(M(\mathbf a))>\varphi^{4T}\)。
反之当 \(\mathbf a=(1,\dots,1)\) 时 \(M(\mathbf a)=(RL)^T\) 为对称正定矩阵，故 \(\sigma_{\max}(M(\mathbf a))=\rho(M(\mathbf a))=\varphi^{2T}\)，从而取到等号。证毕。
\leanverified{goldenRatio\_cube}
\leanverified{goldenRatio\_pow\_five}
\leanverified{goldenRatio\_pow\_six}
\leanverified{paper\_golden\_ratio\_fibonacci\_powers}
\end{proof}

\begin{corollary}[双曲平移长度的线性下界与唯一极小轨]\label{cor:conclusion-event-hyperbolic-translation-golden-minimal}
\leanverified{paper\_conclusion\_event\_hyperbolic\_translation\_golden\_minimal}
在定理 \ref{thm:conclusion-event-ellipse-golden-minimal-stretch} 的设定下，令
\(\ell_{\mathbb H}(g)\) 表示 \(g\in \mathrm{PSL}_2(\RR)\) 的双曲平移长度（即 \(\inf_{z\in\mathbb H}d_{\mathbb H}(z,g\cdot z)\)）。
则对任意 \(T\ge 1\) 与任意 \(\mathbf a\in(\NN_{\ge 1})^T\) 有
\[
\boxed{\;
\ell_{\mathbb H}\!\bigl(M(\mathbf a)\bigr)\ \ge\ 4T\log\varphi\;},
\]
并且等号成立当且仅当 \(\mathbf a=(1,\dots,1)\)。
\end{corollary}
\begin{proof}
对任意 \(g\in \mathrm{SL}_2(\RR)\)，其平移长度满足
\(
\ell_{\mathbb H}(g)=2\log\rho(g)
\)
当 \(g\) 为双曲元时成立；而本处 \(M(\mathbf a)\) 为正整矩阵且 \(\mathrm{tr}\,M(\mathbf a)>2\)，故为双曲元。
由定理 \ref{thm:conclusion-event-ellipse-golden-minimal-stretch} 的证明已得 \(\rho(M(\mathbf a))\ge \varphi^{2T}\)，于是
\(
\ell_{\mathbb H}(M(\mathbf a))\ge 2\log(\varphi^{2T})=4T\log\varphi
\)。
等号情形由谱半径严格单调性与 \(\mathbf a=(1,\dots,1)\) 的对称正定性同理得到。证毕。
\end{proof}
